\documentclass[a4paper, 12pt]{article}

\usepackage[utf8]{inputenc}
\usepackage[T1]{fontenc}
\usepackage{textcomp}
\usepackage{amssymb}
\usepackage{newtxtext} \usepackage{newtxmath}
\usepackage{amsmath, amssymb}
\newtheorem{problem}{Problem}
\newtheorem{example}{Example}
\newtheorem{lemma}{Lemma}
\newtheorem{theorem}{Theorem}
\newtheorem{problem}{Problem}
\newtheorem{example}{Example} \newtheorem{definition}{Definition}
\newtheorem{lemma}{Lemma}
\newtheorem{theorem}{Theorem}
\DeclareMathAlphabet{\mathcal}{OMS}{cmsy}{m}{n}
\usepackage{parskip}

\begin{document}


$\S$ Hace cierto tiempo, soñé que me encontraba en una mesa rodeado de gente que
aprecio y que me aprecia. Entonces, una colega científica decía a los
comensales: «En verdad, $S$ es un joven inteligente»; otro agregó luego
palabras que eran un residuo diurno: «Indeed, he is very accomplished»; otra me
contemplaba con amor y decía lo noble que yo era. Entonces presentí mi pequeñez
y, poniendo las manos sobre la cruz que pende de mi cuello, lloré con amargura.
Nadie comprendió mi llanto, pero les dije: «¡Si tan solo hubiera hecho todo lo
bueno que hice en mi vida, mas sin que nadie supiera que fui yo el que lo
hizo!».

Así batallaba yo con mi vieja arrogancia. Porque solía sentirme mejor que otros
por tener hinchada la cabeza, e ignoraba que la virtud está en el corazón del
hombre. Y cuanto más engordaba mi cabeza, más se torcía mi corazón. Adoraba que
me elogien, y cultivaba ciertas obras por deseo de ser reconocido y celebrado.
Hoy, quisiera nunca más ser elogiado. Cada obra de amor quisiera hacerla en
secreto. El amor es su propio fin.


$\S$ Solía gozar de repetir, en mi ignorancia, aquella sentencia del
\textit{Apocalipsis} que me enseñó mi padre; y, considerándome en cada posición
tomada firme y fuerte, decía «Dios vomita a los tibios». No sólo todo lo juzgaba
sino que me enorgullecía de hacerlo, considerándolo una virtud, porque se me
había enseñado que para ser justo hay que juzgar. Ignoraba la enseñanza
verdadera: \textit{Do not judge, or you too will be judged}. Porque no es justo
el que juzga, sino el que perdona y el que olvida. 

Y, lo que es peor, como un adicto consciente de que está encadenado a un mal, me
decía a mí mismo en secreto: «¡cuánto quisiera no juzgar a nadie!». Veía que, al
desear amar, el juicio me lo impedía: no podía ver con transparencia ni, a su
vez, dejar que nadie me viera. Y odiaba la soledad que yo mismo creaba para mí.

$\S$ Y además, en la cima de la noche de mi vida, me alejé de quienes más
quería, y de las pocas personas que acepté en mi corazón. A mis tres amigos más
hermosos los dejé en mi ciudad natal; a $P$, mi amor, la muerte me la arrebató
de entre las manos; mi familia estaba rota y no tenía a nadie. Por eso a mis
veintiún años consideré la muerte, y seriamente, y sopesé con sumo cuidado cada
factor relevante, y decidí cómo me quitaría la vida. Pero en la tarde señalada
presentí una imagen, que no vi con los sentidos, sino que fue como un latido
interior: un pastor rodeado de luz cuidaba de una oveja, y nada más. Y al sentir
este latido me inundó la paz del mundo. Inmediatamente hablé con un conocido
cristiano, y me dijo «fue Cristo». Y yo le dije: «no fue nadie, fui yo, y lo que
importa no es quién fue, sino que ahora no quiero morir». 

$\S$ Entonces me reconcilié con aquellos a quienes había alejado de mí, y
organicé mi vida, y orienté mis energías a la ciencia y el estudio, y me sentí
sano y dinámico. Pero, aunque ya no estuve solo, seguía lleno de arrogancia, y
cuanto más estudiaba, y cuanto más prosperaba en el trabajo científico, más me
envilecía.

$\S$ Y una amante me confesó un día que al conocerme presintió en mí una
tristeza y una angustia inconmensurables. Y me sorprendió, no por no
reconocerlas, sino por pensar que las ocultaba perfectamente. Y tuvo razón, y
porque me vio desnudo de este modo no quise que me viera nunca más. Porque quien
me ve es quien puede herirme, o así razonaba entonces, sabiéndolo sólo a medias.
Y me alejé sin decirle nada. Por lo que, bajo el aparente dinamismo de mi
reencarnación, seguía siendo oscuro y escurridizo.

$\S$ Recuerdo ahora cierto tiempo feliz, tal vez el más feliz, aunque por mucho
el más colmado de angustias y problemas. Tenía quince años y cursaba el
penúltimo año del secundario, del que egresé con diecisiete, y pasaba mis horas
leyendo en el vasto jardín de mi casa quinta, rodeado de libros y el sonido de
los pájaros. Había encontrado cierta alegría en el abandono del mundo digital
---que, aunque ya denso, no era nada de lo que ahora es---; había guardado mi
computadora en un placard y escribía con una vieja máquina de escribir que mi
padre me consiguió. Él todavía se ríe cariñosamente al recordarme tan jovencito,
sentado a la sombra de lapachos y ceibos, con una pequeña mesa plegable donde
ponía todos mis chiches, interrumpiendo la exuberante orquesta de los trópicos
con la percusión de mi escritura. Estaba trabajando en un libro de poemas que
llamaba las \textit{Flores del odio}, donde sublimaba, con esquivo y dudoso
talento, los retorcidos odios del conflicto familiar. En efecto, aunque un
pintor hubiera visto en ese chico un escenario bucólico y pacífico, yo, por otro
lado, colmaba las páginas de una angustia sin nombre y un amargo resentimiento.
No conoía el perdón, ni me lo habían enseñado; no conocía la piedad, ni la había
sentido nombrar. ¿No es cierto que ese mismo joven, aparentemente tranquilo,
dio muerte sin razón a una hermosa lechuza? Maté para entender que nada es peor
que la muerte; porque, cuando la vi caer, en su elegante majestuosidad, me dije:
«esa criatura es igual a vos, y vos le has dado muerte». Y me odié a mí mismo.

\pagebreak
$\S$ A los catorce años maté, sin razón aparente, a una hermosa lechuza.
Cuando la vi caer, en su elegante sencillez, me dije: «esa criatura es igual a
vos, y vos le has dado muerte». Y me odié a mí mismo. Unos años después pensé en
esa lechuza, y en su extinguida facultad del vuelo, mientras por primera vez
consideraba entregarme a la muerte. Vivía bajo el absoluto convencimiento de que
la totalidad de mi vida había sido una farsa; que cuanto tuvo de hermoso fue una
ilusión y había provisto la diminuta luz de una estrella, y cuanto tenía de
sórdido era terriblemente verdadero y gravitaba con el peso de una grotesca
luna. Subí a la terraza con la resolución de arrojarme al vacío, pero esta
terraza daba al baldío de un regimiento militar, y recordé cierta ocasión de mi
niñez en que mi padre me hizo ver cómo un grupo de soldados excavaba
fatigadamente la tierra. Me explicó que, muchos años atrás, un grupo de hombres
había exterminado y desaparecido a otros hombres, y que ahora estos soldados
buscaban a los desaparecidos. Y me dije a mí mismo: «soy, y he de ser». Porque
sentí que una muerte voluntaria ofendería a aquellos que tuvieron que dar su
vida involuntariamente. Y me dio vergüenza haber matado cualquier criatura,
porque ahora pensaba que la muerte era lo único permanente que había obrado en
esta tierra. Lloraba y me decía: «nunca serás feliz». Porque ignoraba que una luz
brillaba en las tinieblas, aunque las tinieblas no la pudieran comprender.

\pagebreak

$\S$ El sentimiento de vacío no consiste en no sentir nada. Tiene un
\textit{qualia} determinado, pero cuya determinación es precisamente la no
determinación, la vacuidad. Pretendiendo llenarme, busqué el amor:
pero en mi búsqueda quería ser visto, y no ver. Si fui querido, yo sólo presentí
la amorfa intuición de lo que no es. Pero, precisamente por no dar, me vaciaba a
mí mismo. Era como un frasco lleno de humo que se destapa de repente llenando su
alrededor de tinieblas. Era movido por una sed idiota: buscaba dar caza a una
entidad desconocida, una que, cuando finalmente es alcanzada, aunque \textit{en
cierto modo} muere, crea una copia de sí misma en otro lugar desconocido. Y todo
debe comenzar de nuevo.

Ahora sé que sólo ama el que vive en el reino de los cielos. Desde allí uno
desea ver, y no ser visto. Desde allí uno desea dejar su propia luz arder ante
la luz de los demás. Y el camino fue sencillo: no juzgar, perdonarlo todo, y ver
el reino de los cielos en cada ser humano. En los ojos de todas las personas,
descubrir el principio de la vida. Y abandonar mi vieja arrogancia, porque no
soy mejor que nadie y nadie es mejor que yo. 


$\S$ Y, ahora me digo, ¿no son una mentira estas angustias? Se obran en mí como
una ilusión es obrada por un mago habilidoso ante los ojos crédulos de un necio.
He besado una boca y fui feliz---luego, el pensamiento: «this was
meaningless»---luego una extraña paz---luego un vacío---luego ternura. Me engaño
porque imagino que la felicidad es un estado ideal y perfecto; pero, en todas
las cosas, me asombran pensamientos oscuros, infelices o, por lo menos, faltos
de hermosura. Me asombran y me digo: «¿Lo ves? No eres feliz». Como si la
felicidad tuviera como condición necesaria que ningún pensamiento imperfecto se
gestara en nosotros. Pero no es así: se gestan, aparecen, y es nuestro deber
dejarlos ser: verlos pasar como estrellas fugaces en un cielo baldío. La
felicidad es vivir en el reino de los cielos; vivimos en el reino de los cielos
cuando obramos con amor y compasión. Nada importa si un pensamiento oscuro nos
asalta: los pensamientos son como luciérnagas que se prenden y se apagan en la
oscuridad de la consciencia.


























\end{document}



